\chapter{Conclusiones y trabajo futuro}\label{cap:conclusiones}

\section{Conclusiones}

Este trabajo ha planteado la verificación automática de afirmaciones
como un problema estadístico de clasificación condicional sobre pares
\emph{claim}--evidencia, alineado con el dataset FEVER y con la
literatura de inferencia de lenguaje natural. Las contribuciones
principales son:

\begin{enumerate}[leftmargin=*,itemsep=0.3em]
  \item Una formulación explícita del problema y de su evaluación
        (capítulos~\ref{cap:marco-teorico} y~\ref{cap:diseno}).
  \item La implementación y evaluación honesta de un baseline lineal
        TF-IDF + regresión logística multinomial regularizada
        (capítulo~\ref{cap:baseline}).
  \item El \emph{fine-tuning} y comparación de DeBERTa-v3-small como
        modelo de referencia neuronal (capítulo~\ref{cap:deberta}).
  \item El diagnóstico sistemático de calibración
        (capítulo~\ref{cap:calibracion}), identificando que las
        probabilidades del modelo mejoran al baseline pero aún no son
        plenamente calibradas.
  \item La aplicación de contrastes estadísticos formales (McNemar,
        \emph{bootstrap}) para comparar modelos
        (capítulo~\ref{cap:tests}).
  \item Una regla de agregación probabilística que extiende el
        veredicto FEVER con la categoría \textsc{conflicting},
        documentada formalmente en el capítulo~\ref{cap:agregacion} e
        implementada en producción.
\end{enumerate}

\section{Trabajo futuro}

\begin{itemize}[leftmargin=*,itemsep=0.2em]
  \item \textbf{Modelos más grandes.} Evaluar DeBERTa-v3-large o
        modelos NLI de la familia LLaMA / Mistral con \emph{adapters},
        manteniendo las mismas pruebas estadísticas.
  \item \textbf{Recuperación entrenable.} Sustituir la búsqueda web por
        un recuperador denso (DPR, ColBERT) y medir el efecto sobre
        precisión y cobertura.
  \item \textbf{Evaluación cruzada de dominio.} Aplicar el sistema
        entrenado en FEVER a datasets como MultiFC o claims extraídos
        de noticias reales en español, midiendo degradación.
  \item \textbf{Calibración condicional.} Estudiar calibración por clase
        y por longitud de evidencia.
  \item \textbf{Modelado bayesiano del agregador.} Sustituir la regla
        determinista por un modelo gráfico que estime conjuntamente la
        verdad latente y la fiabilidad de cada evidencia.
\end{itemize}
